% =====================================================================
\chapter{C2: handover between TUAVs}\label{ch:c2}
% =====================================================================
\begin{keybox}[colframe=hoc,colback=hoc!5]
A disaster zone needs several TUAVs, and UAVs fly between their areas (``cells''). Neighbouring TUAVs share a link key $K_{ab}$. The old TUAV gives a
moving member a sealed, one-time \textbf{ticket}; the new TUAV admits it after checking the ticket, \textbf{the member's signature} and one fresh
Diffie--Hellman exchange, and replies with \textbf{its own signature}. No TA, no satellite, no full re-login.
\end{keybox}

\section{The gap in the base paper}
\begin{claimbox}
IS4: ``the TUAV, with the effective range covering the entire region, will periodically generate the random number $r_{tu}$ \dots''
\end{claimbox}
The scheme assumes a single TUAV for the whole region. A tethered drone covers a few kilometres at most, so a real deployment has several TUAVs, and the
paper has no TUAV-to-TUAV link, no handover and no leave message. A UAV that crosses into another cell can only start again: US1--US3, GA1 to the TA over
the satellite (three pairings at the TA), and a rebuilt $\Lambda^*(x)$ at the new TUAV -- while the old TUAV is never told that the UAV left.


\section{Two ways to move}
\textbf{Path 1 -- full join (always available).} Any TUAV can verify any credential with $\Ppub$ alone (C1), so a UAV can simply run UA1--UA5 with the
new TUAV. The new TUAV rekeys on arrival and the old TUAV rekeys after the departure.\\
\textbf{Path 2 -- ticket handover (C2).} For a planned move, the member carries a ticket from the old TUAV and does not wait for the new TUAV's next beacon.
The ticket idea goes back to Kerberos~\cite{kerberos94}; what is new here is binding it to signatures in both directions so that a \emph{captured} TUAV
cannot abuse it.

\section{The protocol HO0--HO4}
\begin{stepbox}{hoc}{HO0--HO4 \;Inter-TUAV link and ticket handover}
\steplabel{hoc}{HO0.} Neighbouring TUAVs exchange signed ephemeral values and derive $K_{ab}=\KDF(e_{ab}E_{ba}\cat\ID_a\cat\ID_b\cat\ts)$.\\[2pt]
\steplabel{hoc}{HO1.} Before member $\M$ (pairwise key $K_\M$) leaves, $\Ta$ derives a handover key $\hk=\KDF(K_\M\cat\texttt{HO}\cat\ID_b\cat\ts_h)$ and seals
\[
\Tk=\Enc\big(K_{ab},\ \PID\cat\ep\cat\pk\cat\hk\cat N\cat\mathit{exp};\ \ID_a\cat\ID_b\big)
\]
with a fresh nonce $N$ and an expiry time $\mathit{exp}$; it sends $\Enc(K_\M,\Tk)$ to $\M$.\\[2pt]
\steplabel{uavc}{HO2.} In the new cell $\M$ derives $\hk$ itself, picks $e$, sets $E=eP$ and sends
$\Tk$, $E$, $\MAC(\hk,\Tk\cat E\cat\ts)$ and \textbf{its own signature} on $(\ID_a\cat\ID_b\cat H(\Tk)\cat E)$.\\[2pt]
\steplabel{hoc}{HO3.} $\Tb$ opens $\Tk$ with $K_{ab}$ and rejects if it has expired, if $N$ was used before, if $\M$ has been evicted, if the MAC fails, or if the
signature does not verify under the $\pk$ inside the ticket. It picks $e_b$, derives $K'=\KDF(e_bE\cat\hk\cat\PID\cat\ID_b)$, rekeys its cell, and replies
with $E_b$, $\Enc(K',\GK_b\cat\nu)$ and \textbf{its own signature} on $(\ID_a\cat\ID_b\cat H(\Tk)\cat E\cat E_b)$.\\[2pt]
\steplabel{uavc}{HO4.} $\M$ checks $\Tb$'s signature, derives $K'=\KDF(eE_b\cat\hk\cat\PID\cat\ID_b)$ and opens $\GK_b$. $\Ta$ rekeys once $\M$ has left.
\end{stepbox}

\begin{figure}[H]
\centering
\resizebox{\linewidth}{!}{%
\begin{tikzpicture}[font=\sffamily\small]
  \node[ent=hoc] (a) at (0,0)  {old TUAV $\Ta$};
  \node[ent=uavc] (m) at (6.2,0) {member $\M$};
  \node[ent=hoc] (b) at (12.4,0) {new TUAV $\Tb$};
  \foreach \n in {a,m,b}{ \draw[life] (\n.south) -- ++(0,-11.0); }
  \draw[{Stealth}-{Stealth},hoc,very thick] (0,-0.95) -- node[lbl,above]{HO0 \;signed ECDH $\Rightarrow K_{ab}$} (12.4,-0.95);
  \node[act=hoc,anchor=north] at (0,-1.3) {HO1: $\hk=\KDF(K_\M\cat\texttt{HO}\cat\ID_b\cat\ts_h)$\\ $\Tk=\Enc(K_{ab},\PID\cat\ep\cat\pk\cat\hk\cat N\cat\mathit{exp})$};
  \draw[msg=hoc] (0,-2.7) -- node[lbl,above]{$\Enc(K_\M,\Tk)$} (6.2,-2.7);
  \node[act=uavc,anchor=north] at (6.2,-3.0) {derive $\hk$; pick $e$, $E=eP$\\ \textbf{sign} $(\ID_a\cat\ID_b\cat H(\Tk)\cat E)$};
  \draw[msg=uavc] (6.2,-4.4) -- node[lbl,above]{HO2 \;$\Tk,\ E,\ \MAC(\hk,\Tk\cat E\cat\ts),\ \text{signature}$} (12.4,-4.4);
  \node[act=hoc,anchor=north] at (12.4,-4.7) {HO3: open $\Tk$; check $\mathit{exp}$, $N$ unused,\\ not evicted, MAC, \textbf{member's signature};\\ $K'=\KDF(e_bE\cat\hk\cat\PID\cat\ID_b)$;\\ rekey the cell; \textbf{sign} $(\dots\cat E\cat E_b)$};
  \draw[msg=hoc] (12.4,-7.0) -- node[lbl,above]{$E_b,\ \Enc(K',\GK_b\cat\nu),\ \text{signature of }\Tb$} (6.2,-7.0);
  \node[act=uavc,anchor=north] at (6.2,-7.3) {HO4: \textbf{check $\Tb$'s signature};\\ $K'=\KDF(e\,E_b\cat\dots)$; open $\GK_b$};
  \node[act=hoc,anchor=north] at (0,-8.6) {rekey after $\M$ has left};
  \node[draw=good,fill=good!10,rounded corners,font=\scriptsize\sffamily\bfseries,text=good] at (6.2,-10.3) {$\Tb$ never learns $K_\M$ \;$\cdot$\; one fresh ECDH \;$\cdot$\; signatures both ways};
\end{tikzpicture}}
\caption{Message flow of the C2 handover between the old TUAV $\Ta$, the moving member $\M$ and the new TUAV $\Tb$. Time runs from top to bottom. The two signatures written in bold (by the member in HO2 and by the new TUAV in HO3) are the fixes A2 and A7 of Chapter~\ref{ch:hard}.}
\label{fig:c2flow}
\end{figure}

\section{Why each part is there}
\begin{itemize}[leftmargin=*]
\item \textbf{$\hk$ instead of $K_\M$:} the new TUAV never learns the member's old key, so a later compromise of $\Tb$ does not expose the old cell.
\item \textbf{Fresh ECDH:} $K'$ depends on $e_bE$, so learning $K_{ab}$ or $\hk$ later does not reveal it (forward secrecy).
\item \textbf{Nonce, expiry and both names:} a ticket works once, only before $\mathit{exp}$, and only between these two TUAVs.
\item \textbf{The member's signature (A2):} without it, a captured old TUAV could mint tickets on its own; now a ticket alone admits nobody.
\item \textbf{The new TUAV's signature (A7):} a captured old TUAV knows $\hk$ and could otherwise pose as the new TUAV; now it cannot.
\item \textbf{Rekey on arrival and after departure:} the newcomer cannot read the new cell's past; the leaver cannot read the old cell's future.
\end{itemize}

\section{Security}
We modelled HO0--HO4 in the applied pi calculus and checked it with the ProVerif verifier~\cite{blanchet16} against a Dolev--Yao attacker~\cite{dolevyao}
who also controls a \emph{captured old TUAV} (it knows $K_{ab}$, $K_\M$ and $\hk$).
\begin{center}
\small
\begin{tabularx}{\linewidth}{@{}p{5.6cm} X@{}}
\toprule
\textbf{Model (old TUAV captured)} & \textbf{Result of formal verification}\\
\midrule
Ticket alone admits the member & \no\ attack: a member is ``admitted'' that never started a handover, and the new cell's key leaks \;$\Rightarrow$ \textbf{A2}\\
Member signs, the new TUAV's reply is unsigned & \no\ attack: the member completes a handover with the captured TUAV posing as $\Tb$ \;$\Rightarrow$ \textbf{A7}\\
Both signatures (final protocol) & \yes\ all five properties hold\\
\bottomrule
\end{tabularx}
\end{center}

\begin{theorem}[Handover security, symbolic model]\label{thm:ho}
In the Dolev--Yao model with a compromised old TUAV, the final protocol HO0--HO4 satisfies: (1) injective agreement of $\Tb$ with $\M$ (every accepted
handover was started by $\M$, once); (2) injective agreement of $\M$ with $\Tb$; (3) secrecy of $K'$; (4) secrecy of the new cell key $\GK_b$; and
(5) forward secrecy of $K'$ when $K_{ab}$ and $\hk$ leak after the session.
\end{theorem}
\begin{proof}[Justification]
Established by automatic verification of the model above. Intuitively: (1) needs the member's signature, which the captured TUAV cannot produce
(Theorem~\ref{thm:unforge}); (2) needs $\Tb$'s signature on $E$ and $E_b$; (3)--(5) follow because $K'$ contains $e_bE$, a CDH value unknown to the attacker
even when it holds $K_{ab}$ and $\hk$.
\end{proof}

\section{Cost and latency}
In computation a ticket handover costs about the same as a full join: a few scalar multiplications on each side (one Diffie--Hellman exchange, one
signature and one verification each) plus the cell rekey. Its gain is \emph{latency}: in a full join the member must wait for the new TUAV's next signed
beacon, on average half a beacon period ($500$\,ms for a 1\,s period), whereas with a ticket it starts HO2 at once. On the air, a crossing adds one
member signature (65 bytes) and one signed reply (about 131 bytes). No message ever goes to the TA.

\begin{limitbox}
\begin{itemize}[leftmargin=*]
\item A ticket must never be accepted from a TUAV that has been evicted (Chapter~\ref{ch:c3}): its tickets are sealed under keys a compromised TUAV holds.
Members leaving an evicted TUAV use the full join.
\item C2 protects the handover itself; a compromised \emph{new} TUAV is handled by C3 and C4, not by C2.
\end{itemize}
\end{limitbox}
